\begin{Answer}{8.4}
Nobody in their right mind will choose fruit if cake and ice cream are
available, so there are $3+8=11$ choices.
Just kidding.  There are really $3+8+5=16$ different choices.
\end{Answer}
\begin{Answer}{8.7}
There are 26 choices for each of the first three characters,
and 10 choices for each of the final three characters.
Therefore, there are $26^3\cdot10^3$ possible license plates.
\end{Answer}
\begin{Answer}{8.10}
Every divisor of $n$ is of the form $p_1^{b_1}p_2^{b_2}\cdots p_k^{b_k}$,
where $0\leq b_1\leq a_1$, $0\leq b_2\leq a_2$, $\ldots$, $0\leq b_k\leq a_k$.
(We could also write this as $0\leq b_i\leq a_i$ for $0\leq i\leq k$.)
Therefore there are $a_1+1$ choices for $b_1$,
$a_2+1$ choices for $b_2$, all the way through
$a_k+1$ choices for $b_k$.
Since each of the $b_i$s are independent of each other,
the product rule tells us that the number of divisors of $n$ is
$(a_1 + 1)(a_2 + 1)\cdots (a_k + 1)$.
\end{Answer}
\begin{Answer}{8.11}
Unless the $p_i$ are distinct, the $b_i$s are not independent of each other.
In other words, if the $p_i$s are distinct, then each different choice of the $b_i$s
will produce a different number.  But this is not the case if the $p_i$s are not
distinct.  For instance, if we write $32=2^32^2$, we can get the factor $4$ as
$2^22^0$, $2^12^1$, or $2^02^2$.  Clearly we would count $4$ three times and
would obtain the incorrect number of divisors.
\end{Answer}
\begin{Answer}{8.13}
Write $n = \underbrace{1 + 1 + \cdots + 1}_{n-1 \ +\mathrm{'s}}$.
There are two choices for each plus sign--leave it or perform the addition.
Each of the $2^{n-1}$ ways of making choices leads to a different
expression, and every expression can be constructed this way.
Therefore, there are $2^{n-1}$ such ways of expressing $n$.
\end{Answer}
\begin{Answer}{8.15}
This combines the product and sum rules.
We now have $10+26=36$ choices for each character, and there are $5$ characters,
so the answer is $36^5$.
\end{Answer}
\begin{Answer}{8.16}
Each bit can be either $0$ or $1$, so there are $2^n$ bit strings of length $n$.
\end{Answer}
\begin{Answer}{8.18}
$53\cdot 63^2$;
$53\cdot 63^3$;
$53\cdot 63^{k-1}$.
\end{Answer}
\begin{Answer}{8.21}
It contains at least one repeated digit.  The wording of your answer is very
important.  Your answer should not be ``it has some digit twice'' since this
is vague--do you mean `exactly twice'?  If so, that is incorrect.  If you mean
`at least twice', then it is better to be explicit and say it that way or just say
`repeated'.  To be clear, we don't know that it contains any digit exactly twice,
and we also don't know how many unique digits the number has--it might be $22222222222$,
but it also might be $98765432101$.
\end{Answer}
\begin{Answer}{8.24}
If all
the magenta, all the yellow, all the white, $14$ of the red and $14$
of the blue marbles are drawn, then in among these $8 + 10 + 12 + 14
+ 14 = 58$ there are no $15$ marbles of the same color. Thus we need
$59$ marbles in order to insure that there will be $15$ marbles of
the same color.
\end{Answer}
\begin{Answer}{8.25}
She knows that you are the $25$th person in line.  If everyone gets $4$ tickets,
she will get none, but you will get the $4$ you want.  She can get one or more
tickets if one or more people in front of her, including you, get less than $4$.
\end{Answer}
\begin{Answer}{8.28}
There are seven possible sums, each one a number in
$\{-3,-2,-1,0,1,2,3\}$. By the Pigeonhole Principle, two of the
eight sums must add up to the same number.
\end{Answer}
\begin{Answer}{8.31}
We have $ \lceil{\frac{16}{5}}\rceil = 4$, so some cat has at least four kittens.
\end{Answer}
\begin{Answer}{8.32}
\begin{pfevalenum}
\item
This proof is incomplete.  It kind of argues it for $5$, not $n$ in general.
Even then, the proof is neither clear not complete.  For instance, what are
the 4 `slots'?
\item
They only prove it for $n=2$.  It needs to be proven for {\em any} $n$.
\item
You can't assume somebody had shaken hands with everyone else without some justification.
You certainly can't assume it was any particular person (i.e. person $n$).
Similarly, you can't assume the next person has shaken $n-2$ hands without justifying it.
The final statement is weird (what does `fulfills the contradiction' mean?)
and needs justification (why is it a problem that the last person shakes no hands?).
\end{pfevalenum}
\end{Answer}
\begin{Answer}{8.33}
Notice that if someone shakes $n-1$ hands, then nobody shakes $0$ hands and vice-verse.
Thus, we have two cases.
If someone shakes $n-1$ hands, then the $n$ people can shake
hands with between $1$ and $n-1$ other people.
If nobody shakes hands with $n-1$ people, then the $n$ people can shake
hands with between $0$ and $n-2$ other people. In either case, there are $n-1$
possibilities for the number of hands that the $n$ people can shake.
The pigeonhole principle implies that two people shake hands with the same number of people.

Note:  You cannot say that the two cases are that someone shakes hands with $n-1$ or
someone shakes hands with $0$. It may be that {\em neither} of these is true.
The two cases are someone shakes hands with $n-1$ others or nobody does.  Alternatively,
you could say someone shakes hands with $0$ others or nobody does.
\end{Answer}
\begin{Answer}{8.34}
Choose a particular person of the group, say Charlie. He
corresponds with sixteen others. By the pigeonhole principle,
Charlie must write to at least six of the people about one topic, say
topic I. If any pair of these six people corresponds about topic I,
then Charlie and this pair do the trick, and we are done.
Otherwise, these six correspond amongst themselves only on topics
II or III. Choose a particular person from this group of six, say
Eric. By the Pigeonhole Principle, there must be three of the five
remaining that correspond with Eric about one of the topics, say
topic II. If amongst these three there is a pair that corresponds
with each other on topic II, then Eric and this pair correspond on
topic II, and we are done. Otherwise, these three people only
correspond with one another on topic III, and we are done again.
\end{Answer}
\begin{Answer}{8.38}
$EAT$, $ETA$, $ATE$, $AET$, $TAE$, and $TEA$.
\end{Answer}
\begin{Answer}{8.41}
Since there are $15$ letters and none of them repeat, there are $15!$ permutations
of the letters in the word {\sc uncopyrightable}.
\end{Answer}
\begin{Answer}{8.43}
(a) $5\cdot 7\cdot 6 \cdot 5 \cdot 4 \cdot 3 \cdot 2 = 25,200$.
(b) We condition on the last digit.
If the last digit were $1$ or $5$ then we would have $5$
choices for the first digit and $2$ for the last digit.
Then there are $6$ left to choose from for the second, $5$ for the third, etc.
So this leads to
$$5\cdot 6 \cdot 5\cdot 4 \cdot 3\cdot 2 \cdot 2 = 7,200$$
possible phone numbers.
If the last digit were either $3$ or $7$, then
we would have $4$ choices for the first digit and $2$ for the last.
The rest of the digits have the same number of possibilities as above, so we would
have
$$4\cdot 6\cdot 5 \cdot 4 \cdot 3\cdot 2 \cdot 2  = 5,760   $$
possible phone numbers. Thus the total number of phone numbers is
$$7200 + 5760 = 12,960.$$
\end{Answer}
\begin{Answer}{8.45}
Label the letters $T_1$, $A_1$, $L_1$, and $L_2$.  There are $4!$ permutations of these letters.
However, every permutation that has $L_1$ before $L_2$ is actually identical to one having
$L_1$ before $L_2$, so we have double-counted.  Therefore, there are $4!/2=12$ permutations
of the letters in $TALL$.
\end{Answer}
\begin{Answer}{8.46}
$TALL$, $TLAL$, $TLLA$, $ATLL$, $ALTL$, $ALLT$, $LLAT$, $LALT$, $LATL$, $LLTA$, $LTLA$, and $LTAL$.
That makes $12$ permutations, which is exactly what we said it should be in Exercise~\ref{exer:tallperm1}.
\end{Answer}
\begin{Answer}{8.47}
Following similar logic to the previous few examples, since we have one letter that is repeated
three times, and a total of $5$ letters, the answer is $5!/3!=20$.
\end{Answer}
\begin{Answer}{8.48}
Ten of them are
$AIEEE$, $AEIEE$, $AEEIE$, $AEEEI$, $EAIEE$, $EAEIE$, $EAEEI$, $EEAIE$, $EEAEI$, $EEEAI$.
The other ten are identical to these, but with the $A$ and $I$ swapped.
\end{Answer}
\begin{Answer}{8.51}
We can consider $SMITH$ as one block along with the remaining
$5$ letters $A$, $L$, $G$, $O$, and $R$. Thus, we are
permuting $6$ `letters', all of which are unique.  So there are
$6!=720$ possible permutations.
\end{Answer}
\begin{Answer}{8.54}
(a) $5\cdot 8^6 =1310720$.
(b) $5\cdot 8^5 \cdot 4 = 655360 $.
(c) $5\cdot 8^5 \cdot 4 = 655360 $.
\end{Answer}
\begin{Answer}{8.58}
(a) $\dfrac{7\cdot 6\cdot 5 \cdot 4\cdot 3}{1\cdot 2\cdot 3\cdot 4 \cdot 5} = 21$.
(b) $\dfrac{12\cdot 11}{1\cdot 2} = 66$.
(c) $\dfrac{10\cdot 9\cdot 8\cdot 7\cdot 6}{1\cdot 2\cdot 3\cdot 4 \cdot 5} = 252$.

\noindent
(d) $\dfrac{200\cdot 199\cdot 198\cdot 197}{1\cdot 2\cdot 3\cdot 4} =64,684,950$.
(e) $1$.
\end{Answer}
\begin{Answer}{8.61}
(a) $\displaystyle\binom{17}{15} = \displaystyle\binom{17}{2}=\dfrac{17\cdot 16}{1\cdot 2} = 136$.
(b) $\displaystyle\binom{12}{10} = \displaystyle\binom{12}{2}=\dfrac{12\cdot 11}{1\cdot 2} = 66$.

\noindent
%(c) $\displaystyle\binom{10}{6}=\displaystyle\binom{10}{4}=\dfrac{10\cdot 9\cdot 8\cdot 7}{1\cdot 2\cdot 3\cdot 4} = 210$.
(c) $\displaystyle\binom{200}{196}=\displaystyle\binom{200}{4}=\dfrac{200\cdot 199\cdot 198\cdot 197}{1\cdot 2\cdot 3\cdot 4} =64,684,950$.
%\noindent
(d) $\displaystyle\binom{67}{66}=\displaystyle\binom{67}{1}=67/1=67$.
\end{Answer}
\begin{Answer}{8.65}
$12, 13, 14, 15, 23, 24, 25, 34, 35, 45$.
\end{Answer}
\begin{Answer}{8.68}
\begin{solevalenum}
\item
This solution does not take into account which woman was selected
and which 15 of the original 16 are left, so this is not correct.
\item
This solution has two problems.  First, it counts things multiple times.  For instance,
any selection that contains both Sally and Kim will be counted twice--once
when Sally is the first woman selected and again when Kim is selected first.
Second, the product rule should have been used instead of the sum rule.  Of course, that
hardly matters since it would have been wrong anyway.
\item
This solution is correct.
\end{solevalenum}
\end{Answer}
\begin{Answer}{8.70}
To count the  number of shortest routes from $A$ to $B$
that pass through point $O$, we count the number of
paths from $A$ to $O$ (of which there are $\binom{5}{3} = 10$) and
the number of paths from $O$ to $B$ (of which there are
$\binom{4}{3} = 4$). Using the product rule,
the desired number of paths is
$\displaystyle\binom{5}{3}\binom{4}{3} = 10\cdot 4 = 40$.
\end{Answer}
\begin{Answer}{8.71}
\begin{solevalenum}
\item
This answer is incorrect since it will count some of the committees multiple times.
If you did not come up with an example of something that gets counted multiple times,
you should do so to convince yourself that this answer is incorrect.
\item
This solution is incorrect since it does not take into account which man and woman
were selected  and which 14 of the original 16 are left.
\end{solevalenum}
\end{Answer}
\begin{Answer}{8.72}
There are $\binom{16}{5}$ possible committees.  Of these, $\binom{9}{5}$  contain only men and
$\binom{7}{5}$ contain only women.
Clearly these two sets of committees do not overlap.
Therefore, the number of committees that contain at least one man and at least one woman is
$\binom{16}{5}-\binom{9}{5} -\binom{7}{5}$.
\end{Answer}
\begin{Answer}{8.73}
Because we subtracted the size of both of these from the total number of possible
committees.  If the sets intersected, we would have subtracted some possibilities twice
and the answer would have been incorrect.
\end{Answer}
\begin{Answer}{8.75}
\begin{solevalenum}
\item
This solution is incorrect since it double counts some of the possibilities.
\item
This solution is incorrect because it does not take into account the requirement
that one course from each group must be taken.
\end{solevalenum}
\end{Answer}
\begin{Answer}{8.76}
\begin{solevalenum}
\item
This solution is incorrect since it counts some of the possibilities multiple times.
\item
This solution is incorrect because it does not take into account the requirement
that one course from each group must be taken.
\end{solevalenum}
\end{Answer}
\begin{Answer}{8.79}
Using $10$ bars to separate the meat and $3$ stars to represent the slices,
we can see that this is exactly the same as the previous two examples.  Thus,
the solution is $\binom{13}{10}=\binom{13}{3}=286.$
\end{Answer}
\begin{Answer}{8.84}
\begin{eqnarray*}
%\! is a small negative space. This is too wide, but barely and that fixes it.
(2x - y^2)^4 &\!\! = & \!\!\!\! \binom{4}{0}(2x)^4 + \binom{4}{1}(2x)^3(-y^2) + \binom{4}{2}(2x)^2(-y^2)^2 + \binom{4}{3}(2x)(-y^2)^3 + \binom{4}{4}(-y^2)^4 \\
 & = & (2x)^4 + 4(2x)^3(-y^2) + 6(2x)^2(-y^2)^2 + 4(2x)(-y^2)^3 + (-y^2)^4 \\
& = & 16x^4 - 32x^3y^2 + 24x^2y^4 - 8xy^6 + y^8
\end{eqnarray*}
\end{Answer}
\begin{Answer}{8.85}
$$
\begin{array}{lll}
(\sqrt{3} + \sqrt{5})^4 & = & (\sqrt{3})^4 +
4(\sqrt{3})^3(\sqrt{5})  + 6(\sqrt{3})^2(\sqrt{5})^2
+ 4(\sqrt{3})(\sqrt{5})^3 + (\sqrt{5})^4 \\
& = & 9 + 12\sqrt{15} + 90 +  20\sqrt{15} + 25 \\
& = & 124 + 32\sqrt{15}
\end{array}
$$
\end{Answer}
\begin{Answer}{8.87}
Using a little algebra and the binomial theorem,  we can see that
$$\dsum_{k=1}^{n} \binom{n}{k}3^k
=\dsum_{k=0}^{n} \binom{n}{k}3^{k}-1
=\dsum_{k=0}^{n} \binom{n}{k}1^{n-k}3^{k}-1
=(1+3)^{n}-1=4^n-1.$$
\end{Answer}
\begin{Answer}{8.91}
Let $A$ be the set of camels eating wheat and $B$ be the set of camels eating barley.
We know that $|A|=46$, $|B|=57$, and $|A\cup B|=100-10=90$.
We want $|A\cap B|$.  By Theorem~\ref{thm:two_set_incl_excl} (solving it for $|A\cap B|$),
$$|A\cap B| = |A|+|B|-|A \cup B| = 46+57-90=13.$$
\end{Answer}
\begin{Answer}{8.95}
Using Theorem~\ref{thm:three_set_incl_excl}, we know that
$28+29+19-14-10-12+8=48\%$  watch at least one of these sports.
That leaves $52\%$ that don't watch any of them.
\end{Answer}
\begin{Answer}{8.96}
Let $C$ denote the set of people
who like candy, $I$ the set of people who like ice cream, and $K$
denote the set of people who like cake. We are given that $|C|
= 816$, $|I| = 723$, $|K| = 645$, $|C\cap I| = 562$,
$|C\cap K| = 463$, $|I\cap K| = 470$, and $|C\cap I\cap K| = 310$. By Inclusion-Exclusion we have
$$\begin{array}{lll}
|C \cup I \cup K| & = & |C| + |I| + |K| \\ & & -
|C \cap I| - |C \cap K |
- |I \cap C| \\ & &  + |C \cap I \cap K| \\
& = & 816 + 723 + 645 - 562 - 463 - 470 + 310 \\ & = & 999.
\end{array}$$
The investigator miscounted, or probably did not report one person
who may not have liked any of the three things.
\end{Answer}
\begin{Answer}{8.98}
We can either use inclusion-exclusion for four sets or use a few applications
of inclusion-exclusion for two sets.  Let's try the latter.

Let $A$ denote the set of
those who lost an eye, $B$ denote those who lost an ear, $C$ denote
those who lost an arm and $D$ denote those losing a leg. Suppose
there are $n$ combatants.
 Then
$$\begin{array}{lll}
n & \geq & | A \cup B| \\ & = &  | A| + | B| - | A \cap B| \\ & = & .7n + .75n - | A\cap B|,
\end{array}$$
$$\begin{array}{lll}
n  &\geq & | C \cup D| \\ & = & | C| + | D| - | C \cap D| \\ & = & .8n + .85n - |C\cap D|.
\end{array}$$
This gives
$$| A \cap B| \geq .45n,$$
$$| C \cap D| \geq .65n.$$This means that
$$\begin{array}{lll}n &\geq & | (A\cap B) \cup (C \cap D)| \\ & = & | A \cap B| + | C \cap D| - | A \cap B \cap C \cap D|\\
& \geq &  .45n + .65n - | A \cap B \cap C \cap D|,\end{array}$$
whence
$$| A \cap B \cap C \cap D| \geq .45 + .65n - n = .1n.$$This means
that at least $10\%$ of the combatants lost all four members.
\end{Answer}
